\section{The Datum, the Approximants and the Phase Clock}
\label{sec:datum}

We define the datum $f$ of Theorem~\ref{thm:main} and its approximants
$f^{(n)}$; the solutions $\omega^{(n)}$ with these data form the
$2$-admissible sequence to which Corollary~I.\ref{I-cor:defect} is applied in
the proof of Proposition~\ref{prop:reduction}. We then split each $\omega^{(n)}$
into its rings, describe each ring in its own units by its perturbation and by
the field that the other rings exert on it, and define the phase clock, which
fixes, from $\omega^{(n)}$ itself, the time at which each ring becomes active and
the time at which its unstable coefficient is expected to reach the size
$\delta_0$.

\begin{definition}
\label{def:datum}
Let $A_1>0$, $\mu\in(1,2]$ and $\sbar\in(0,a/16]$, and define $A_k$, $s_k$,
$r_k$ ($k\ge1$) and $f$ by Equation~\eqref{eq:datum}.
\end{definition}

\begin{proposition}
\label{prop:integrability}
In Definition~\ref{def:datum}, the sets $\{ar_k\le|x|\le r_k\}$, $k\ge1$, are
pairwise disjoint.
The function $f$ is radial and nonnegative, $f\in (L^1\cap L^q)(\R^2)$ for every
$q\in[1,\infty)$, and $f\notin L^\infty(\R^2)$. Moreover
$\sum_{k\ge1}A_k^qr_k^2<\infty$ for every $q\in[1,\infty)$.
\end{proposition}
\begin{proof}
Since $\mu>1$ and $\sbar\le a/16$, $s_k\le\sbar/2<a$, so $r_{k+1}<ar_k$ and every
annulus of index $i>k$ lies in the open disk $B_{ar_k}$, because $r_i\le r_{k+1}$;
this disk does not meet the annulus of index $k$, so the annuli are pairwise
disjoint. Moreover
$r_k=\prod_{j<k}s_j=\sbar^{k-1}(2\mu)^{-k(k-1)/2}$.

By Lemma~I.\ref{I-lem:ring-facts}, $\Phi$ is radial and continuous, $0\le\Phi<1$,
$\Phi\not\equiv0$ and $\supp\Phi\subset\{a\le|y|\le1\}$. So the $k$-th term of
the series is radial, nonnegative and supported in the $k$-th annulus; since the
annuli are disjoint, at every point at most one term is nonzero, $f$ is radial and nonnegative, and, by
monotone convergence and the formula for $r_k$, for $q\in[1,\infty)$,
\begin{gather*}
\int f^q\dd x=\sum_{k\ge1}A_k^qr_k^2\|\Phi\|_{L^q}^q,\qquad
\|\Phi\|_{L^q}^q\le|\{a\le|y|\le1\}|\le\pi,\\
A_k^qr_k^2=A_1^q\mu^{q(k-1)}\sbar^{2(k-1)}(2\mu)^{-k(k-1)}=A_1^q\big(\mu^q\sbar^2(2\mu)^{-k}\big)^{k-1}.
\end{gather*}
If $(2\mu)^k\ge2\mu^q$, which holds for all large $k$ because $2\mu>1$, the
bracket is at most $\frac12$, because $\sbar<1$, so
$A_k^qr_k^2\le A_1^q2^{1-k}$ for all large $k$. Hence $\sum_kA_k^qr_k^2<\infty$
and $f\in L^q$. Finally, $f\ge A_k\Phi(\cdot/r_k)$, and, $\Phi$ being continuous,
$A_k\Phi(\cdot/r_k)>\frac12A_k\max\Phi>0$ on a nonempty open set; since $\mu>1$
gives $A_k\to\infty$, $f\notin L^\infty$.
\end{proof}

The approximants keep the first $n$ rings of $f$ and replace ring $n$, the
innermost of them, whose time scale $1/A_n$ is the shortest, by the seeded ring
$\Phi^{\kseed}$ of Lemma~\ref{lem:seed}. This is a nonnegative, $3$-fold
rotationally symmetric function, equimeasurable with $\Phi$, whose coefficient
$\ellst(0)(\Phi^{\kseed}-\Phi)$ along the unstable eigenfunction $\phist(0)$ has
modulus between $\kseed/2$ and $2\kseed$. The relay described in the
Introduction starts from this seed.

\begin{definition}
\label{def:approx}
Assume \textup{(N1)}--\textup{(N3)}, and let $\kL$ and $\Phi^\kappa$ be as in
Lemma~\ref{lem:seed}. For an integer $n\ge2$ and $\kseed\in(0,\kL]$ put
$f^{(n)}=\sum_{k<n}A_k\Phi(\cdot/r_k)+A_n\Phi^{\kseed}(\cdot/r_n)$, let $\omega^{(n)}$ be the solution of Theorem~\ref{thm:smooth} with datum
$f^{(n)}$, let $v^{(n)}=\KBS\ast\omega^{(n)}$, and let $\flow_t$ be the flow of
$v^{(n)}$.
\end{definition}

\begin{lemma}
\label{lem:adm}
In Definition~\ref{def:approx}, $f^{(n)}\in C_c^\infty(\R^2)$,
$f^{(n)}$ and $\omega^{(n)}(t)$ are $3$-fold rotationally symmetric for every
$t$, and, for fixed $\kseed$, $f^{(n)}\to f$ in $(L^1\cap L^q)(\R^2)$ for every
$q\in[1,\infty)$ as $n\to\infty$. Moreover $(f^{(n)},\omega^{(n)})_{n\ge2}$ is a
$2$-admissible sequence for $f$ on $[0,\infty)$ in the sense of
Definition~I.\ref{I-def:adm} (the sequence starts at
$n=2$; the shift $n\mapsto n+1$ gives the indexing of Definition~I.\ref{I-def:adm}).
\end{lemma}
\begin{proof}
By Lemmas~I.\ref{I-lem:ring-facts} and \ref{lem:seed}, $\Phi$ and $\Phi^{\kseed}$ lie
in $C_c^\infty(\R^2)$, are $3$-fold rotationally symmetric ($\Phi$ is radial)
and equimeasurable, so
they have the same $L^q$ norms. Hence $f^{(n)}\in C_c^\infty(\R^2)$ is
$3$-fold rotationally symmetric, and every $\omega^{(n)}(t)$ is $3$-fold
rotationally symmetric, by Theorem~\ref{thm:smooth}\textup{(6)} with the rotation $\Orot$. Let
$q\in[1,\infty)$. In
$f-f^{(n)}=A_n(\Phi-\Phi^{\kseed})(\cdot/r_n)+\sum_{k>n}A_k\Phi(\cdot/r_k)$, the
first term has $L^q$ norm at most $2A_nr_n^{2/q}\|\Phi\|_{L^q}$, by the triangle
inequality and the equality of the $L^q$ norms of $\Phi$ and $\Phi^{\kseed}$; the
terms of the sum have disjoint supports, by Proposition~\ref{prop:integrability}, so
\[
\|f-f^{(n)}\|_{L^q}\le\Big(2A_nr_n^{2/q}
+\Big(\sum_{k>n}A_k^qr_k^2\Big)^{1/q}\Big)\|\Phi\|_{L^q},
\]
and both terms on the right tend to $0$ as $n\to\infty$ because
$\sum_kA_k^qr_k^2<\infty$ (Proposition~\ref{prop:integrability}): $A_n^qr_n^2$ is
a term, and $\sum_{k>n}A_k^qr_k^2$ a tail, of this convergent series. Finally, by
Theorem~\ref{thm:smooth}\textup{(7)}, for $q=2$ and every
$T>0$ the restriction of $\omega^{(n)}$ to $[0,T)$ is a renormalized weak
solution in $X_2(T)$ with datum $f^{(n)}$; so $\omega^{(n)}$ is a renormalized
weak solution on $[0,\infty)$, and with the convergence just proved this is the
$2$-admissibility of Definition~I.\ref{I-def:adm}.
\end{proof}

Each ring is studied in its own scale, in which the $k$-th term of $f^{(n)}$
becomes $\Phi$, or $\Phi^{\kseed}$ if $k=n$, so that the results of
Section~\ref{sec:single}, stated in ring units, can be applied to it. For $A,r>0$ the function
$\tilde\omega(\tau,y)=A^{-1}\omega^{(n)}(\tau/A,ry)$ is the solution of
Theorem~\ref{thm:smooth} with datum $A^{-1}f^{(n)}(r\,\cdot)$, with velocity
$(\KBS\ast\tilde\omega)(\tau,y)=(Ar)^{-1}v^{(n)}(\tau/A,ry)$ and flow
$y\mapsto r^{-1}\flow_{\tau/A}(ry)$. Indeed, $\KBS$ is homogeneous of degree
$-1$, which gives the velocity; then
$\partial_\tau\tilde\omega+(\KBS\ast\tilde\omega)\cdot\nabla\tilde\omega
=A^{-2}(\partial_t\omega^{(n)}+v^{(n)}\cdot\nabla\omega^{(n)})(\tau/A,ry)=0$, and
Theorem~\ref{thm:smooth}\textup{(5)} identifies $\tilde\omega$ with that solution.
With $A=A_k$ and $r=r_k$, $\tilde\omega$ is $\omega^{(n)}$ in \textbf{ring-$k$
units}: points $x$ become $y=x/r_k$, vorticities are divided by $A_k$, and
$\tau=A_kt$ is the \textbf{ring-$k$ time}. For a velocity field $V$ we write
$V_{[k]}(y)=V(r_ky)/(A_kr_k)$ for $V$ in ring-$k$ units.

The next definition splits $\omega^{(n)}$ into its rings and names the
quantities that the later sections estimate. The set $S_k(t)$ carries ring $k$
along the flow, and $\omega_{\ge k}$ is the part of $\omega^{(n)}$ formed by ring
$k$ and the rings inside it. In ring-$k$ units, the perturbation $w_k$ is the
deviation of ring $k$ from $\Phi$; $2\pi G_k$ is the circulation of the point
vortex by which the linearization $\Lin(G_k)$ of Theorem~I.\ref{I-thm:linear}
represents the rings inside ring $k$; and the external field $h_k$ is the field
of all the other rings on $N'$ minus the field of this point vortex, that is,
the field $h$ of Lemma~\ref{lem:reduction}.

The number $\varepsilon_k$ is the small factor in the bounds of
Lemma~\ref{lem:transfer}\textup{(1)}--\textup{(3)} on the field, minus its
monopole, that the material inside ring $k$ induces on ring $k$. The
normalized moment $q_{k+1}$ is the mode-$3$ moment through which this material
drives the unstable coefficient of ring $k$ (Lemma~\ref{lem:coupling}), and the
Duhamel majorant $\maj_k$ bounds, up to a constant factor, the $H^{s_h}$ norm
of $w_k$ while ring $k$ is quiet and $\maj_k$ is small
(Theorem~\ref{thm:quiet}\textup{(c)}).

\begin{definition}
\label{def:regions}
Fix $n\ge2$ and the objects of Definition~\ref{def:approx}. For $1\le k\le n$:
\begin{itemize}
\item $f^{(n)}_k$ is the $k$-th term of $f^{(n)}$, that is,
  $f^{(n)}_k=A_k\Phi(\cdot/r_k)$ for $k<n$ and
  $f^{(n)}_n=A_n\Phi^{\kseed}(\cdot/r_n)$;
  \[
  S_k=\supp f^{(n)}_k,\quad S_k(t)=\flow_t(S_k),\quad D_k(t)=\bigcup_{i=k}^nS_i(t),\quad
  D_{n+1}(t)=\emptyset,
  \]
  and $\omega_{\ge k}(t)=\omega^{(n)}(t)\mathbf 1_{D_k(t)}$;
\item $N_k=r_kN$, $N'_k=r_kN'$, $\Nflat_k=r_k\Nflat$, $\NN_k=r_k\NN$,
  $\NNm_k=r_k\NNm$ (the annuli of Section~I.\ref{I-sec:ring-def} and of Section~\ref{sec:II-prelim} dilated by $r_k$);
\item while $S_k(t)\subset N_k$, the \textbf{perturbation} of ring $k$ at the
  time $t$, in ring-$k$ units, is
  $w_k(t,y)=A_k^{-1}\omega^{(n)}(t,r_ky)\mathbf 1_{S_k(t)}(r_ky)-\Phi(y)$;
\item the \textbf{inner circulation} $G_k$ is defined by
  $2\pi G_kA_kr_k^2=\int\omega_{\ge k+1}(0)$; we put
  $\varepsilon_k=\mu s_k^{m+2}$, and, if $G_k\le\Gone$, $\sigma_k=\sigma(G_k)$,
  $\lambda_k=\lamst(G_k)$, $x_k=\sigma_k/\mu$ and
  $\beta_k=\beta(G_k)s_k^{m+2}/(2\pi)$;
\item for $k<n$, the \textbf{normalized moment} of the inner region is
  \[
  q_{k+1}(t)=\overline{\Qm[\omega_{\ge k+1}(t)]}/(A_{k+1}r_{k+1}^{m+2})
  \]
  (the support of $\omega_{\ge k+1}(t)$ is compact and does not contain $0$, by
  Lemma~\ref{lem:regions-basic}\textup{(2)} below, so $\Qm$ applies);
\item at a time $t$ with $D_{k+1}(t)\subset B_{ar_k/8}$, the \textbf{inner field}
  of ring $k$ at ring-$k$ time $\tau=A_kt$ is
  \[
  \hin_k(\tau)=\KBS\ast\nu^{\mathrm{in}}_k-G_k\frac{y^\perp}{|y|^2}\quad\text{on }N',\qquad
  \nu^{\mathrm{in}}_k(y)=A_k^{-1}\omega_{\ge k+1}(t,r_ky);
  \]
  at a time $t$ with $\big(S_1(t)\cup\dots\cup S_{k-1}(t)\big)\cap B_{6r_k}=\emptyset$
  (the rings $i<k$ lie outside $B_{6r_k}$), the \textbf{outer field} is
  $\hout_k(\tau)=\KBS\ast\nu^{\mathrm{out}}_k$ on $N'$, with
  $\nu^{\mathrm{out}}_k(y)=A_k^{-1}\omega^{(n)}(t,r_ky)\mathbf 1_{S_1(t)\cup\dots\cup S_{k-1}(t)}(r_ky)$;
  when both are defined, $h_k=\hin_k+\hout_k$ is the \textbf{external field} of
  ring $k$;
\item for an integer $M_h\ge1$ (from Section~\ref{sec:relay} on, $M_h$ is the
  integer of \textup{(SA2)} in Definition~\ref{def:standing}) we put
  $\etin_k(\tau)=\|\hin_k(\tau)\|_{C^{M_h}(N')}$,
  $\etout_k(\tau)=\|\hout_k(\tau)\|_{C^{M_h}(N')}$ and $\eta_k=\etin_k+\etout_k$,
  as functions of ring-$k$ time;
\item if $G_k\le\Gone$ and $\eta_k$ is defined on $[0,\tau]$ (it is then
  continuous there, as we show after Lemma~\ref{lem:regions-basic}), the
  \textbf{Duhamel majorant} of ring $k$ is
  \[
  \maj_k(\tau)=\int_0^\tau e^{\sigma_k(\tau-v)}\eta_k(v)\dd v ,
  \]
  and $\maj^{\mathrm{in}}_k$ is defined in the same way from $\etin_k$.
\end{itemize}
\end{definition}

\begin{lemma}
\label{lem:regions-basic}
In Definition~\ref{def:regions} the following hold for every $t\ge0$.
\begin{enumerate}[label=\textup{(\arabic*)},beginpenalty=10000]
\item The sets $S_1,\dots,S_n$ are compact, pairwise disjoint and invariant
  under $\Orot$, and none contains $0$. Moreover
  $S_k\subset\{ar_k\le|x|\le r_k\}$ for $k<n$, $S_n\subset\NNm_n$, and
  $f^{(n)}\mathbf 1_{S_k}=f^{(n)}_k$.
\item $\flow_t\circ\Orot=\Orot\circ\flow_t$ and $\flow_t(0)=0$. The sets
  $S_1(t),\dots,S_n(t)$ are compact, pairwise disjoint and invariant under
  $\Orot$, none contains $0$, and
  $\supp\omega^{(n)}(t)\subset S_1(t)\cup\dots\cup S_n(t)$.
\item For $1\le k\le n$,
  \begin{equation}
  \label{eq:omega-ge}
  \omega^{(n)}(t)\mathbf 1_{S_k(t)}=f^{(n)}_k\circ\flow_t^{-1},\qquad
  \omega_{\ge k}(t)=\Big(\sum_{i=k}^nf^{(n)}_i\Big)\circ\flow_t^{-1}
  =\sum_{i=k}^n\omega^{(n)}(t)\mathbf 1_{S_i(t)},
  \end{equation}
  and $\int\omega^{(n)}(t)\mathbf 1_{S_k(t)}=A_kr_k^2\|\Phi\|_{L^1}$. These
  functions lie in $C_c^\infty(\R^2)$, are nonnegative and are $3$-fold
  rotationally symmetric.
\item If $S_k(t)\subset N_k$, then $w_k(t)\in C_c^\infty(\R^2)$ is real valued,
  $3$-fold rotationally symmetric and supported in $N$; so $w_k(t)\in X^s$ for every integer
  $s\ge0$. Moreover $\int w_k(t)=0$, $\|w_k(t)\|_{L^1}\le2\|\Phi\|_{L^1}$, and
  $\|w_k(t)\|_{L^1}\le|N|^{1/2}\|w_k(t)\|_{L^2}$.
\end{enumerate}
\end{lemma}

\begin{proof}
(1) For $k<n$, $f^{(n)}_k$ is radial and supported in $\{ar_k\le|x|\le r_k\}$
(Lemma~I.\ref{I-lem:ring-facts}); $f^{(n)}_n$ is $3$-fold rotationally symmetric and supported in
$\NNm_n$ (Lemma~\ref{lem:seed}). So each $S_k$ is compact, $\Orot$-invariant and
avoids $0$. Let $i<k$. By these support properties, $S_k\subset\overline B_{5r_k/4}$, and
$\frac54r_k\le\frac54s_ir_i<ar_i$, because $r_k\le r_{i+1}=s_ir_i$ and $s_i\le\sbar/2\le a/32$, while
$S_i\subset\{|x|\ge ar_i\}$; so $S_k\cap S_i=\emptyset$, and
$f^{(n)}\mathbf 1_{S_k}=f^{(n)}_k$.

(2) By Lemma~\ref{lem:adm}, $\omega^{(n)}(t)\circ\Orot=\omega^{(n)}(t)$; with
$\KBS(\Orot x)=\Orot\KBS(x)$ and the change of variables $y=\Orot y'$ this gives
$v^{(n)}(t,\Orot x)=\Orot v^{(n)}(t,x)$. So $t\mapsto\Orot^{-1}\flow_t(\Orot y)$
solves the equation of $\flow_t(y)$ with the same initial value; since
$v^{(n)}$ is $C^1$ and $\nabla v^{(n)}$ is bounded on $[0,t]\times\R^2$
(Theorem~\ref{thm:smooth}\textup{(2)}), Gronwall's inequality gives uniqueness,
and $\flow_t\circ\Orot=\Orot\circ\flow_t$. At $x=0$, $v^{(n)}(t,0)=\Orot
v^{(n)}(t,0)$, so $v^{(n)}(t,0)=0$ and $\flow_t(0)=0$. As $\flow_t$ is a
homeomorphism (Theorem~\ref{thm:smooth}\textup{(4)}), the sets
$S_k(t)=\flow_t(S_k)$ inherit compactness, disjointness and $\Orot$-invariance
from (1), $0\notin S_k(t)$, and
$\supp\omega^{(n)}(t)=\flow_t(\supp f^{(n)})\subset\bigcup_kS_k(t)$, because
$\omega^{(n)}(t)=f^{(n)}\circ\flow_t^{-1}$.

(3) Since $\omega^{(n)}(t)=f^{(n)}\circ\flow_t^{-1}$ and
$\mathbf 1_{S_k(t)}=\mathbf 1_{S_k}\circ\flow_t^{-1}$, (1) gives the first
identity, and (2), by which $D_k(t)$ is the disjoint union of
$S_k(t),\dots,S_n(t)$, gives the other two. Each $f^{(n)}_k$ lies in
$C_c^\infty(\R^2)$, is nonnegative and $3$-fold rotationally symmetric (Lemmas~I.\ref{I-lem:ring-facts}
and \ref{lem:seed}), and $\flow_t^{-1}$ is
a $C^\infty$ diffeomorphism commuting with $\Orot$; this gives the regularity
and symmetry. Since $\det D\flow_t=1$,
$\int\omega^{(n)}(t)\mathbf 1_{S_k(t)}=\int f^{(n)}_k=A_kr_k^2\int\Phi
=A_kr_k^2\|\Phi\|_{L^1}$, as $\Phi\ge0$; for $k=n$ we use that the nonnegative equimeasurable
functions $\Phi^{\kseed}$ and $\Phi$ have the same integral.

(4) The two terms of
$w_k(t,y)=A_k^{-1}(\omega^{(n)}(t)\mathbf 1_{S_k(t)})(r_ky)-\Phi(y)$ are real,
lie in $C_c^\infty(\R^2)$, are $3$-fold rotationally symmetric (by (3) and
Lemma~I.\ref{I-lem:ring-facts}) and are supported in $N$, because
$S_k(t)\subset N_k$ and $\supp\Phi\subset\{a\le|y|\le1\}\subset N$
(Lemma~\ref{lem:annuli}). By (3) and the change of variables $x=r_ky$, both are
nonnegative with integral, hence $L^1$ norm, $\|\Phi\|_{L^1}$; so $\int w_k(t)=0$ and $\|w_k(t)\|_{L^1}\le2\|\Phi\|_{L^1}$.
The last inequality is the Cauchy-Schwarz inequality on $N$.
\end{proof}

The fields of Definition~\ref{def:regions} depend continuously on time: if
$J\subset[0,\infty)$ is a compact interval with $D_{k+1}(t)\subset B_{ar_k/8}$
for $t\in J$, then $\tau\mapsto\hin_k(\tau)$ is continuous from $A_kJ$ into
$C^M(N')$ for every integer $M\ge0$. Indeed, $\hin_n=-G_ny^\perp/|y|^2$ is
constant, because $D_{n+1}(t)=\emptyset$. Let $k<n$ and $g=\sum_{i>k}f^{(n)}_i$, supported in the compact set
$D_{k+1}(0)$. By Equation~\eqref{eq:omega-ge}, $\omega_{\ge k+1}(t)=g\circ\flow_t^{-1}$,
so the change of variables $y=\flow_t(y')$, with $\det D\flow_t=1$, gives
$(\KBS\ast\omega_{\ge k+1}(t))(x)=\int\KBS(x-\flow_t(y'))g(y')\dd y'$. As
$(t,y')\mapsto\flow_t(y')$ is continuous (Theorem~\ref{thm:smooth}\textup{(4)}),
$\rho=\max\{|\flow_t(y')|:t\in J,\ y'\in D_{k+1}(0)\}$ exists, and $\rho<\frac a8r_k$,
because the maximum is attained at a point of one of the sets
$\flow_t(D_{k+1}(0))=D_{k+1}(t)\subset B_{ar_k/8}$.
Let $U$ be the set of the $x$ at distance less than $\frac a{16}r_k$ from $N'_k$;
as the inner radius of $N'_k$ is $\frac a4r_k$, every $x\in U$ has
$|x|>\frac3{16}ar_k$. So, for $x\in U$, $t\in J$ and $y'\in D_{k+1}(0)$,
$|x-\flow_t(y')|\ge\frac3{16}ar_k-\rho>\frac a{16}r_k$. Since
$|\partial^\beta\KBS(z)|\le C_\beta|z|^{-1-|\beta|}$, the derivatives
$\partial^\beta_x$, $|\beta|\le M$, of the integrand are bounded, for $x\in U$
and $t\in J$, by a constant times $|g(y')|$, which is integrable. So we may
differentiate under the integral sign, and each of these derivatives of the field is
continuous on $J\times U$ by dominated convergence, hence uniformly continuous on
the compact set $J\times N'_k$; since $G_ky^\perp/|y|^2$ does not depend on
$\tau$, passing to ring-$k$ units gives the claim. In the same way, if the rings $i<k$ lie outside $B_{6r_k}$
for $t\in J$, then $\tau\mapsto\hout_k(\tau)\in C^M(N')$ is continuous on $A_kJ$:
replace $g$ by $\sum_{i<k}f^{(n)}_i$ and $U$ by $B_{4r_k}$, which contains $N'_k$,
and use $|x-\flow_t(y')|\ge2r_k$ for
$x\in B_{4r_k}$ and $y'\in S_1\cup\dots\cup S_{k-1}$, which holds because
$\flow_t(y')$ then lies in $S_1(t)\cup\dots\cup S_{k-1}(t)$, outside $B_{6r_k}$. Consequently $\etin_k$, $\etout_k$ and
$\eta_k$ are continuous on every interval on which they are defined, so the
integrals that define $\maj_k$ and $\maj^{\mathrm{in}}_k$ exist.

If on $[0,T]$ both $D_{k+1}(t)\subset B_{ar_k/8}$ and the rings $i<k$ lie outside
$B_{6r_k}$, then, by Lemma~\ref{lem:regions-basic}, the solution in ring-$k$
units, split into its parts on $S_k(t)$, on $D_{k+1}(t)$ and on
$S_1(t)\cup\dots\cup S_{k-1}(t)$ (rescaled), satisfies the hypotheses of
Lemma~\ref{lem:reduction} on $[0,A_kT]$ with $G=G_k$, and $h_k$ is the field $h$
there. Indeed, by Lemma~\ref{lem:regions-basic}\textup{(2)} and \textup{(3)}, the
three parts have pairwise disjoint supports, are $3$-fold rotationally symmetric
and are transported by the flow of the solution in ring-$k$ units. The part on
$D_{k+1}(t)$ is $\nu^{\mathrm{in}}_k$, whose integral at $t=0$ is $2\pi G_k$ by
the definition of $G_k$, so $G=G_k$. The two conditions on $[0,T]$, read in ring-$k$
units, imply the support conditions of Lemma~\ref{lem:reduction}. While
$S_k(t)\subset N_k$, the function $w=\omega_{\mathrm r}-\Phi$ of
Lemma~\ref{lem:reduction}\textup{(3)} at ring-$k$ time $A_kt$ is $w_k(t)$. This is how Lemma~\ref{lem:reduction} and Section~\ref{sec:single} are
applied to ring $k$.

Part~I and Section~\ref{sec:single} treat the linearization $\Lin(G)$ for
$G\in[0,\Gone]$. The next lemma shows that every inner circulation $G_k$ lies in
this range when $\sbar$ is small, with a bound independent of $n$.

\begin{lemma}
\label{lem:inner-circ}
In Definition~\ref{def:regions}, $\int\omega_{\ge k+1}(t)$ is independent of
$t$, and $0\le G_k\le C_G\mu s_k^2\le C_G\mu\sbar^2$ for $1\le k\le n$, with
a constant $C_G$ depending only on $\Phi$; moreover $G_n=0$. In particular, if
$C_G\mu\sbar^2\le\Gone$, then $G_k\in[0,\Gone]$ for $1\le k\le n$.
\end{lemma}
\begin{proof}
Let $F\colon\R\to\R$ be continuous with $F(0)=0$. Since $\flow_t$ preserves
area (Theorem~\ref{thm:smooth}\textup{(4)}), Equation~\eqref{eq:omega-ge} gives
$\int F(\omega_{\ge k}(t))\dd x=\int F\big(\sum_{i=k}^nf^{(n)}_i\big)\dd x$. The
$f^{(n)}_i$ are supported in the pairwise disjoint sets $S_i$
(Lemma~\ref{lem:regions-basic}\textup{(1)}), so, as $F(0)=0$, the function
$F\big(\sum_{i=k}^nf^{(n)}_i\big)$ is the sum of the $F(f^{(n)}_i)$. Hence, for
every such $F$ and every $1\le k\le n+1$,
\begin{equation}
\label{eq:conserved}
\int F\big(\omega_{\ge k}(t)\big)\dd x=\sum_{i=k}^n\int F\big(f^{(n)}_i\big)\dd x .
\end{equation}
With $F(s)=s$, $\int\omega_{\ge k+1}(t)$ does not depend on $t$, and, as
$\int f^{(n)}_i=A_ir_i^2\|\Phi\|_{L^1}$ (Lemma~\ref{lem:regions-basic}\textup{(3)}
at $t=0$), $2\pi G_kA_kr_k^2=\|\Phi\|_{L^1}\sum_{i=k+1}^nA_ir_i^2$. So $G_k\ge0$
and $G_n=0$, as the sum is then empty. Since $A_{i+1}r_{i+1}^2=\mu s_i^2\,A_ir_i^2$ and
$\mu s_i^2\le2(a/16)^2<\frac12$, each term of the sum is less than half the
previous one, and comparison with a geometric series shows that the sum is at most
$2A_{k+1}r_{k+1}^2=2\mu s_k^2A_kr_k^2$; so
$G_k\le\|\Phi\|_{L^1}\mu s_k^2/\pi\le C_G\mu s_k^2$, and
$s_k\le\sbar$.
\end{proof}

The next lemma bounds the part $\omega_{\ge k}$ of the solution formed by ring
$k$ and the rings inside it, uniformly in $n$ and $t$. Its velocity bound controls the
velocity field of $\omega_{\ge k}$ itself; in the proof of
Proposition~\ref{prop:confine} it is one of the two terms of the speed bound for
the material inside a ring, and its constant $\Cstar$ enters the confinement
radius there.

\begin{lemma}
\label{lem:speed}
In Definition~\ref{def:regions}, there are absolute constants $\Mvort$
and $\Cstar$ such that for all $1\le k\le n$, $t\ge0$ and $q\in\{1,4\}$,
\[
\begin{gathered}
\|\omega_{\ge k}(t)\|_{L^q}\le\Mvort^{1/q}A_kr_k^{2/q},\qquad
\|\KBS\ast\omega_{\ge k}(t)\|_{L^\infty}\le\Cstar A_kr_k,\\
\|\omega_{\ge k}(t)\|_{L^1}=\|\Phi\|_{L^1}\sum_{i=k}^nA_ir_i^2\le2\|\Phi\|_{L^1}A_kr_k^2 .
\end{gathered}
\]
\end{lemma}
\begin{proof}
By Equation~\eqref{eq:conserved} with $F(s)=|s|^q$ and the equimeasurability of
$\Phi^{\kseed}$ and $\Phi$ (Lemma~\ref{lem:seed}),
$\|\omega_{\ge k}(t)\|_{L^q}^q=\|\Phi\|_{L^q}^q\sum_{i=k}^nA_i^qr_i^2$. For
$q\in\{1,4\}$ and $i\ge k$ we have
$A_{i+1}^qr_{i+1}^2=\mu^qs_i^2A_i^qr_i^2\le\mu^4\sbar^2A_i^qr_i^2\le\frac18A_i^qr_i^2$,
because $\mu\le2$ and $\sbar\le a/16<\frac1{16}$ give $\mu^4\sbar^2<\frac1{16}$;
so $\sum_{i=k}^nA_i^qr_i^2\le\frac87A_k^qr_k^2$. With $q=1$ this gives the
identity and the bound for $\|\omega_{\ge k}(t)\|_{L^1}$. As $0\le\Phi<1$ and
$\supp\Phi\subset\{a\le|y|\le1\}$ (Lemma~I.\ref{I-lem:ring-facts}), $\|\Phi\|_{L^q}^q\le\pi$, and
$\|\omega_{\ge k}(t)\|_{L^q}^q\le\frac{8\pi}7A_k^qr_k^2\le\Mvort A_k^qr_k^2$.
Finally, let $g=\omega_{\ge k}(t)\in (L^1\cap L^4)(\R^2)$; we may assume $g\ne0$. Let
$x\in\R^2$ and $R>0$; by Lemma~I.\ref{I-lem:BS-basic}(a), $2\pi|\KBS(x-y)|=|x-y|^{-1}$.
We split the integral of $|g(y)||x-y|^{-1}$ over $y$ into the near field and the
far field.
On the near field, H\"older's inequality with exponents $4,4/3$ and
integration in polar coordinates give
\[
\int_{|x-y|<R}|g(y)||x-y|^{-1}\dd y\lesssim\|g\|_{L^4}R^{1/2}.
\]
The far-field integral is at most $\|g\|_{L^1}R^{-1}$. Choosing
$R=(\|g\|_{L^1}/\|g\|_{L^4})^{2/3}$ balances these bounds and yields
\[
\|\KBS\ast g\|_{L^\infty}
\lesssim\|g\|_{L^1}^{1/3}\|g\|_{L^4}^{2/3}
\lesssim A_kr_k.
\]
The implicit constants are absolute; choose $\Cstar$ to dominate the last one.
\end{proof}

The phase clock turns the relay described in the Introduction into definite
times. It is defined recursively, from the innermost ring outward. Ring $n$
becomes active at time $0$. Ring $k$ becomes active at the time $\tsh_k$, and
$T_k$ is the time at which its unstable coefficient is expected to reach the
size $\delta_0$. Ring $k-1$ then becomes active after a pinning phase
of length $\Tpin$ in ring-$k$ time. At the end of this pinning phase the kick
time of ring $k-1$, defined below, equals $\Tpin$; there
Lemma~\ref{lem:coupling}\textup{(2)} evaluates the kick that ring $k$ and the
rings inside it give to ring $k-1$.

\begin{definition}
\label{def:clock}
Let $\delta_0,\Tpin>0$. In Definition~\ref{def:regions}, assume $G_k\le\Gone$
for $1\le k\le n$. Define recursively, downward in $k$:
\begin{itemize}
\item $\tsh_n=0$;
\item if $\tsh_k$ is defined and $S_k(\tsh_k)\subset N_k$:
  $c_k(0)=\ellst(G_k)\big(w_k(\tsh_k)\big)$ and $\ksh_k=|c_k(0)|$;
\item if moreover $\ksh_k\in(0,\delta_0)$:
  $\hat T_k=\sigma_k^{-1}\log(\delta_0/\ksh_k)$, $T_k=\tsh_k+\hat T_k/A_k$, and, if
  $k\ge2$, $\tsh_{k-1}=T_k+\Tpin/A_k$.
\end{itemize}
The clock is \textbf{defined up to level $k$} if $S_i(\tsh_i)\subset N_i$ and
$\ksh_i\in(0,\delta_0)$ for every $i\in[k,n]$, so that $T_i$ is defined for
these $i$. Put $T_{n+1}=0$, $\tash_k=A_k\tsh_k$, $W_k=[T_{k+1},T_k]$, the
\textbf{active phase} $\Phase_k=[\tsh_k,T_k]$ and the \textbf{pinning phase}
$\Pin_k=[T_{k+1},\tsh_k]$ ($k<n$). Ring $k$ is \textbf{quiet} on $[0,\tsh_k]$,
\textbf{active} on $\Phase_k$, and \textbf{debris} after $T_k$. On $\Phase_k$ we
use the \textbf{active time} $u=A_k(t-\tsh_k)$; if $h_k$ is defined on
$[\tsh_k,t]$, the running coefficient of ring $k$ is, for $0\le u\le A_k(t-\tsh_k)$,
\begin{equation}
\label{eq:running-k}
c_k(u)=c_k(0)+\int_0^ue^{-\lambda_kv}\,\ellst(G_k)\big(F_k(v)\big)\dd v ,
\end{equation}
where $F_k(v)=-h_k(\tash_k+v)\cdot\nabla\Phi$ on $N'$ and $F_k(v)=0$ elsewhere.
The \textbf{kick time} of ring $k<n$ is $\tauK=A_{k+1}(t-T_{k+1})$; it equals
$\Tpin+\mu u$ on $\Phase_k$, and $\tau_0=-A_{k+1}T_{k+1}$ is the kick time of
$t=0$.
\end{definition}

The window $W_k$ is the union of the pinning phase $\Pin_k$ and the active phase
$\Phase_k$ of ring $k$ (for $k=n$, of $\Phase_n$ alone).

The integral in Equation~\eqref{eq:running-k} exists, so $c_k(u)$ is well
defined: its integrand is continuous in $v$, since the map
$v\mapsto h_k(\tash_k+v)\in C^0(N')$ is continuous (as shown after
Lemma~\ref{lem:regions-basic}) and
$|\ellst(G_k)(-h\cdot\nabla\Phi)|\le\Cls\|h\|_{C^0(N')}$ for continuous $h$
(Lemma~I.\ref{I-lem:uniform}).

The relay lemmas of Section~\ref{sec:relay} are causal, and the bootstrap of
Section~\ref{sec:bootstrap}, a continuity argument in time, applies them at a
time $t$ under hypotheses on $[0,t]$. So the phase boundaries that enter them at
time $t$ must be determined by $\omega^{(n)}$ on $[0,t]$, in the following sense, and
Lemma~\ref{lem:clock} below shows that they are. A quantity built from $\omega^{(n)}$ by
Definitions~\ref{def:regions} and \ref{def:clock} is
\textbf{determined by $\omega^{(n)}$ on $[0,\bar t]$} if it keeps its value when
$\omega^{(n)}$ is replaced in these definitions by any
$\tilde\omega\in C^\infty([0,\infty)\times\R^2)$ such that $\tilde\omega=\omega^{(n)}$
on $[0,\bar t]\times\R^2$, every $\tilde\omega(t)$ has compact support, and
$\KBS\ast\tilde\omega$ is smooth with every derivative bounded on $[0,T]\times\R^2$
for every $T$; the flow $\flow$ is then replaced by the flow of
$\KBS\ast\tilde\omega$, and $f^{(n)}$, the sets $S_k$ and the parameters are kept.

\begin{lemma}
\label{lem:clock}
In Definition~\ref{def:clock}, assume that the clock is defined up to level $k$.
\begin{enumerate}[label=\textup{(\arabic*)}]
\item $0=\tsh_n<T_n<\tsh_{n-1}<T_{n-1}<\dots<\tsh_k<T_k$, and $T_k<\tsh_{k-1}$
  if $k\ge2$.
\item For every $i\in[k,n]$, each of $\tsh_i$, $\ksh_i$, $T_i$ and, if $i\ge2$,
  $\tsh_{i-1}$ is determined by $\omega^{(n)}$ on $[0,\tsh_i]$.
\item Hence, for every $t\ge0$, the set $\mathcal J_t$ of the $i\in[k,n]$ with
  $\tsh_i\le t$, and the numbers $\tsh_i$ and $T_i$ for $i\in\mathcal J_t$, are
  determined by $\omega^{(n)}$ on $[0,t]$. In particular, so are the phase
  boundaries $\tsh_i$, $T_i$ ($k\le i\le n$) not exceeding $t$, and the phase of
  each ring $i\in[k,n]$ at time $t$ in the sense of Definition~\ref{def:clock}.
\end{enumerate}
\end{lemma}
\begin{proof}
(1) Let $i\in[k,n]$. By Lemma~\ref{lem:inner-circ}, $G_i\ge0$, and $G_i\le\Gone$ by
the assumption of Definition~\ref{def:clock}; so Proposition~I.\ref{I-prop:SU} gives
$\sigma_i=\sigma(G_i)>\frac{211}{2500}$. Moreover $\ksh_i<\delta_0$, as the clock
is defined up to level $k$, so $\hat T_i>0$
and $T_i>\tsh_i$; if $i\ge2$, then $\tsh_{i-1}-T_i=\Tpin/A_i>0$. These inequalities, for $i=n,n-1,\dots,k$,
give (1).

(2) Let $\bar t\ge\tsh_i$, let $\tilde\omega$ be as above, and mark with a tilde
the objects built from $\tilde\omega$. The velocities $\KBS\ast\tilde\omega$ and
$v^{(n)}$ agree on $[0,\bar t]\times\R^2$ and are Lipschitz in $x$, so, by
uniqueness for the equation of the flow, the two
flows agree on $[0,\bar t]$; and $\tilde G_j=G_j$, because
$\tilde\omega(0)=f^{(n)}$, so the numbers $\sigma_j$ do not change either. We show by downward induction on $j\in[i,n]$ that
$\tilde t^\sharp_j=\tsh_j$, $\tilde\kappa^\sharp_j=\ksh_j$, $\tilde T_j=T_j$, and
$\tilde t^\sharp_{j-1}=\tsh_{j-1}$ if $j\ge2$. First, $\tilde t^\sharp_n=0=\tsh_n$.
If $j\in[i,n]$ and $\tilde t^\sharp_j=\tsh_j$, then $\tsh_j\le\tsh_i\le\bar t$ by (1), so
$\tilde S_j(\tsh_j)=S_j(\tsh_j)\subset N_j$ and $\tilde w_j(\tsh_j)=w_j(\tsh_j)$;
hence $\tilde\kappa^\sharp_j=\ksh_j\in(0,\delta_0)$, and
Definition~\ref{def:clock} gives $\tilde T_j=T_j$ and
$\tilde t^\sharp_{j-1}=\tsh_{j-1}$. This completes the induction and proves (2).

(3) Let $\tilde\omega$ agree with $\omega^{(n)}$ on $[0,t]\times\R^2$, and let
$\tilde{\mathcal J}_t$ be the set of the $i\in[k,n]$ for which
$\tilde t^\sharp_i$ is defined and at most $t$. Let $i_0=\min\mathcal J_t$, which
exists because $\tsh_n=0$. By (1), $\tsh_i$ decreases as $i$ increases, so
$\mathcal J_t=[i_0,n]$, and $\tsh_i,T_i>t$
for $i\notin\mathcal J_t$. By (2) with $\bar t=t$ and $i=i_0$, the numbers
$\tsh_i$, $T_i$, $i\ge i_0$, are the same for $\tilde\omega$, so
$[i_0,n]\subset\tilde{\mathcal J}_t$, and, if $i_0>k$,
$\tilde t^\sharp_{i_0-1}=\tsh_{i_0-1}>t$; the proof of (1) applies to the clock
of $\tilde\omega$ wherever it is defined, so every defined $\tilde t^\sharp_i$ with $i<i_0$ exceeds
$t$. Hence $\tilde{\mathcal J}_t=\mathcal J_t$. These data decide the phase of
ring $i\in[k,n]$ at $t$. If $i\notin\mathcal J_t$, then $t<\tsh_i<T_i$ and ring
$i$ is quiet and neither active nor debris, and for $\tilde\omega$ the number
$\tilde t^\sharp_i$ is either undefined or greater than $t$, so ring $i$ is
neither active nor debris and is quiet at $t$ in the reading of
Definition~\ref{def:loc} in Section~\ref{sec:several}, which takes $t\le\tsh_i$ to
be true when $\tsh_i$ is not defined. If $i\in\mathcal J_t$, ring $i$ is
quiet if and only if $t=\tsh_i$, active if and only if $t\le T_i$, and debris if
and only if $t>T_i$.
\end{proof}
